% Solutions to the problems of Chapter 8 (Open Challenges and Future Directions).
% Included by backmatter/solutions.tex through \InputIfFileExists.
% [October 2026 round] New file (blue).
\begin{addedblock}
\begin{sol}{prob:ch08:1}
(a) For each sensor $i$ and each detection associated to a confirmed track, the monitor computes the innovation $\boldsymbol{\nu}_{i,k}$, its covariance $\mathbf{S}_{i,k}$ and the NIS $d^2_{i,k}$ of \eqref{eq:nis}; over a sliding window of $N$ samples (typically $N=20$--$50$, i.e.\ $2$--$5$~s at $10$~Hz) it keeps the windowed mean NIS $\bar d^2_i$, the windowed mean innovation (a bias estimate), the fraction of expected returns actually received (detection rate), and, when two sensors observe the same track, the Mahalanobis distance between the single-sensor estimates. (b) The windowed mean NIS is compared with $\chi^2_{Nn_y,1-\alpha}/(Nn_y)$, where $n_y$ is the measurement dimension; $\alpha$ is chosen from the tolerated false-alarm rate per window (for instance $10^{-3}$) and a $k$-out-of-$m$ persistence rule further reduces false alarms at the cost of detection delay. The detection-rate threshold is set from nominal statistics, for instance fewer than half of the expected returns during one second. (c) A reliability factor $\alpha_i\in[0,1]$ is derived as a decreasing function of the excess NIS, e.g.\ $\alpha_i=\min\{1,\ \chi^2_{Nn_y,1-\alpha}/(Nn_y\,\bar d^2_i)\}$, and multiplied by a second factor derived from the detection rate; the factor is used both as an evidence-theoretic discount (mass $1-\alpha_i$ is transferred to the frame of discernment) and as a covariance inflation $\mathbf{R}_i/\alpha_i$ in the tracker, so that the two layers of Fig.~\ref{fig:hybrid-pipeline} stay consistent. (d) A blocked sensor shows a collapse of the detection rate while the NIS of the detections that still arrive is normal: the sensor is declared unavailable, the tracker continues with the remaining modalities, and the planner is informed of the reduced field of view and of the larger covariance (a larger protection level). A miscalibrated sensor shows a normal detection rate but a biased innovation whose sign is consistent and, for a temporal offset, grows with the relative velocity (Problem~\ref{prob:ch08:3}): the sensor is discounted, online calibration is triggered, and the planner receives a degraded-integrity flag with the inflated covariance. If either condition persists beyond a pre-defined time and the remaining sensors cannot meet the required integrity, the policy requests a minimal-risk manoeuvre.
\end{sol}

\begin{sol}{prob:ch08:2}
(a) With edges $0,0.2,\dots,1$ the first bin is empty; bin $[0.2,0.4)$ contains two samples (mean confidence $0.285$, accuracy $0.5$, gap $0.215$); bin $[0.4,0.6)$ three samples ($0.517$, $0$, gap $0.517$); bin $[0.6,0.8)$ three samples ($0.717$, $0.667$, gap $0.050$); bin $[0.8,1]$ four samples ($0.900$, $0.750$, gap $0.150$). Hence $\mathrm{ECE}=(2\cdot0.215+3\cdot0.517+3\cdot0.050+4\cdot0.150)/12=0.2275\approx0.23$. (b) The Brier score is $0.2313$; the accuracy is $6/12=0.50$ and the mean confidence $0.656$. (c) The mean confidence exceeds the accuracy, so the detector is over-confident, and the confidences must be flattened towards $0.5$, which requires $T>1$. Minimising the negative log-likelihood (script \repopath{ch08\_\allowbreak p2\_\allowbreak ece.py} in \repopath{companion\_\allowbreak repo/\allowbreak exercises/\allowbreak ch08/\allowbreak solutions/\allowbreak }) gives $T^\ast\approx1.82$, reducing the NLL from $0.656$ to $0.623$ and the ECE from $0.228$ to $0.132$. (d) With two to four samples per bin, the per-bin accuracy has a standard error of about $0.25$--$0.35$, and the absolute value in the ECE makes the estimate biased upwards for small samples. ECE values are only meaningful with thousands of samples per bin, and estimators with adaptive (equal-mass) bins or kernel smoothing are preferable; the exercise illustrates the mechanics, not a practical evaluation.
\end{sol}

\begin{sol}{prob:ch08:3}
(a) A measurement taken at $t_k-\delta$ but interpreted at $t_k$ reports the range $r(t_k-\delta)=r(t_k)+v\delta$ for a target closing at speed $v$, so the innovation is $\nu_k=v\delta+e_k$ with $e_k\sim\mathcal{N}(0,\sigma^2)$ and $b=v\delta$: $b=0.30$~m for $\delta=20$~ms and $b=0.75$~m for $\delta=50$~ms. (b) $\mathbb{E}[\nu_k^2/\sigma^2]=(b^2+\sigma^2)/\sigma^2=1+b^2/\sigma^2$, i.e.\ $2.0$ and $7.25$. (c) At $50$~ms the expected NIS exceeds the gate, so a correct measurement is rejected on average; since $\nu_k/\sigma\sim\mathcal{N}(b/\sigma,1)$, the rejection probability is $P(\lvert Z+b/\sigma\rvert>1.96)$, which equals $0.17$ at $20$~ms and $0.70$ at $50$~ms, against $0.05$ without offset. (d) The windowed test fires on average when $1+b^2/\sigma^2>1.571$, i.e.\ $\lvert b\rvert>0.3\sqrt{0.571}=0.227$~m, giving $\delta_{\min}=0.227/15=15$~ms at $15$~m/s and $113$~ms at $2$~m/s. (e) The same offset is both most harmful and most detectable at motorway speeds; in slow urban traffic a $100$~ms offset is invisible to the test and almost harmless ($0.2$~m). Since $\partial h/\partial\delta=v$, the offset is observable whenever $v\neq0$ and can be appended to the state vector and estimated (augmented-state or unified spatio-temporal calibration), or removed at the source by hardware time-stamping. The script \repopath{ch08\_\allowbreak p3\_\allowbreak timestamp\_\allowbreak offset.py} in the same folder reproduces these numbers by Monte Carlo simulation.
\end{sol}

\begin{sol}{prob:ch08:4}
Feed three trackers with the camera-only, LiDAR-only, and fused detections of the simulation-trained detector and record, over real recordings, (i) the windowed mean NIS per modality, which should be close to one per degree of freedom if the learned covariance transfers (the ratio of real to simulated mean NIS measures how much of the uncertainty model survived); (ii) the distribution of the conflict mass $K$ between the camera and LiDAR branches, which should match its simulated distribution; (iii) temporal consistency, measured by the fraction of detections associated to confirmed tracks and by the flicker rate of track existence; and (iv) the distribution of confidence scores, compared with its simulated counterpart by a two-sample test. Accept a modality if its mean NIS stays within the $\chi^2$ bounds of the window, the shift of the $K$ distribution is below a pre-set threshold, and persistence is not lower than in simulation by more than a margin; otherwise first rescale the covariance (a one-parameter fix that the NIS ratio provides directly) and re-test, and re-train only if the score distribution has shifted as well. A small labelled subset adds what label-free statistics cannot see: false negatives (a missed object produces no innovation), the calibration error of the scores, and a sanity check of the thresholds themselves.
\end{sol}

\begin{sol}{prob:ch08:5}
(a) Object identifier and provenance (sender identity, sensor lineage, and the list or hop count of contributing sources); time stamp in a declared time base with the sender's clock accuracy; reference frame and the sender's own pose with its covariance; the definition of the state vector; the full covariance matrix rather than a scalar confidence; an existence or detection probability usable as a basic belief assignment; a declared reliability or health status; and the motion model used by the sender for prediction. (b) A default covariance that is too small gives the message an excessive weight in the covariance-intersection combination and makes the fused covariance over-optimistic (inconsistent); one that is too large discards the information. Without provenance the same track received directly and through a relaying vehicle is treated as two independent pieces of evidence, the covariance shrinks by a factor of about two, and the estimate becomes confidently wrong, which is exactly the double counting that covariance intersection prevents when the receiver knows that the inputs may be correlated; without provenance the receiver must apply intersection to everything and loses accuracy. (c) Whenever the receiver and the sender observe the same object, the receiver computes the NIS of the sender's track against its own estimate and maintains a running trust score; the score is mapped to a discount factor that inflates the sender's covariance before intersection and discounts its belief masses. Trust must be built over time, corroborated by independent sources, and protected by message authentication so that a spoofer cannot inherit the trust of a legitimate sender.
\end{sol}

\begin{sol}{prob:ch08:6}
Parse the free-text label onto the frame of discernment $\Theta$ (the known classes plus an ``other'' element); an unambiguous label with self-reported confidence $c$ gives $m(\{A\})=c$ and $m(\Theta)=1-c$, and an ambiguous description gives mass to the corresponding subset (for instance ``some animal'' to the set of animal classes), which is precisely what evidence theory allows and a probability vector does not. The discount factor $\alpha$ is estimated on a held-out set as the empirical precision of the model's labels as a function of its reported confidence, so that the discounted assignment $m^\alpha(A)=\alpha\,m(A)$, $m^\alpha(\Theta)=1-\alpha+\alpha\,m(\Theta)$ is calibrated. A large conflict mass $K$ with the explicit sensors means that the semantic source disagrees with physical evidence and should be kept visible, as in the transferable belief model, rather than renormalised away; it indicates either a semantic error of the model or a sensor failure, and which of the two is the case is decided by the consistency monitors of the classical layer. The source should be inserted at the evidence monitor of Fig.~\ref{fig:hybrid-pipeline}, acting on the class part of the state only: it may add hypotheses and increase the class uncertainty, but it must not be able to reduce the existence probability of a track that is consistent with the Kalman tracker's innovation tests, which is enforced by a floor on the mass that the combination may remove from geometrically confirmed objects.
\end{sol}
\end{addedblock}
